% ============================================================
% Figure: Attack signal split (RQ4)
% ============================================================
\begin{figure*}[t]
\centering
\fbox{\parbox{0.95\textwidth}{\centering\todo{Two stacked sub-plots:
(a) commitment\_fail\_rate vs strength for 5 content-tamper
attacks, all rise linearly with strength. (b)
missing\_leaves\_rate vs strength for 4 structural attacks, also
linear. Together they show \memmark{} detects \emph{which}
audit layer caught \emph{which} attack — content vs structural —
a single conflated ``tamper detection rate'' would lose this.}
\\[6pt] \vspace{5cm}}}
\caption{\textbf{Attack signal split — content vs structural
tamper}. Content-tamper attacks (manual\_edits, paraphrase\_rewrite,
edge\_relabel, subgraph\_reanchor, supersession) trigger
\texttt{commitment\_fail\_rate} since leaf bytes change in place.
Structural attacks (pruning, dedup, poisoning, compaction) trigger
\texttt{missing\_leaves\_rate} since the leaf set changes size.
Reporting both lets reviewers see exactly which audit layer caught
which attack.}
\label{fig:attack-signal-split}
\end{figure*}
